\documentclass{article}
\usepackage[T1,T2A]{fontenc}
\usepackage[utf8]{inputenc}
\usepackage[english,russian]{babel}
\usepackage{xcolor}
\usepackage{amsmath}
\usepackage{amsfonts,bm}
\usepackage{tikz}
\usepackage{listings}
\usetikzlibrary{patterns}
\usetikzlibrary{quotes,angles}
\usepackage{relsize}
\usepackage{float}

\title{Решение Задачи №39}
\author{Серёжи Фуканчика\\34 Ю класс}
\date{Февраль 2020}

\begin{document}

\maketitle

\section{Задача}
Решите неравенство
$$f(f(x))\ge{}(f(x))^2$$
где $f(x)=2x^2-1$.

\section{Решение}
Сделаю один шаг замены $f(x)$ на $t$:
$$f(t)\ge{}t^2$$

Раскрою $f(t)$:
$$2t^2-1\ge{}t^2$$

Вычту $t^2$ слева и справа:
$$t^2-1\ge{}0$$

Прибавлю $1$ слева и справа:
$$t^2\ge{}1$$

Снова возвращаюсь к $f(x)$:
$$(f(x))^2\ge{}1$$
$$(2x^2-1)^2\ge{}1$$
$$(2x^2)^2+2(2x^2)(-1)+1\ge{}1$$

$-1$ превращю в минус:
$$(2x^2)^2-2(2x^2)+1\ge{}1$$

вычту $1$ слева и справа:
$$(2x^2)^2-2(2x^2)\ge{}0$$

вынесу $2x^2$ за скобку:
$$(2x^2)((2x^2)-2)\ge{}0$$

можно поделить на $2$, это не повлияет на ответ (\small{почему?}):
$$x^2(x^2-1)\ge{}0$$

$(x^2-1)$ это ``квадрат разности''. Разложу его на множители:
$$x^2(x-1)(x+1)\ge{}0$$

Здесь перемножаются три выражения - $x^2$, $x-1$ и $x+1$. У них разные промежутки где они больше, меньше, либо равны нулю. Эти промежутки частично пересекаются. По условию задачи \textbf{мы эти выражения перемножаем}. Из арифметики мы знаем, что

\begin{tabular}{|c|c|}
\hline
     $+\cdot{}+$ & $+$ \\
     $+\cdot{}-$ & $-$ \\
     $-\cdot{}+$ & $-$ \\
     $-\cdot{}-$ & $+$ \\
\hline
\end{tabular}

т.е. в тех точках где наши две прямые обе больше нуля мы имеем $+\cdot{}+=+$. Если в какой-то точке одна из прямых ниже нуля а вторая - выше у нас будет $+\cdot{}-=-$. И т.п.

Давай я последовательно нарисую (расположу промежутки один над другим, а потом объединю их):

\begin{center}
  \begin{tikzpicture}
    % vertical lines
    \draw[->] (0cm,-4.5) -- (0cm,2cm);
    \draw[dashed] (1,-4.5) -- (1,1);
    \draw[dashed] (-1,-4.5) -- (-1,1);

    % axes
    \draw[->] (-5cm,0cm) -- (5cm,0cm) node[right,fill=white] {$x$};
    \draw[->] (-5cm,-2cm) -- (5cm,-2cm) node[right,fill=white] {$x$};
    \draw[->] (-5cm,-4cm) -- (5cm,-4cm) node[right,fill=white] {$x$};

    % functions
    \draw (10/3,1.1) node {$x^2$};
    \draw (10/3,-1.1) node {$x-1$};
    \draw (10/3,-3.1) node {$x+1$};


    % 0 for all three functions
    \draw (0.2,-0.5) node {$0$};
    \draw (0.2,-2.5) node {$0$};
    \draw (0.2,-4.5) node {$0$};

    % 1 for all three functions
    \draw (1.2,-0.5) node {$1$};
    \draw (1.2,-2.5) node {$1$};
    \draw (1.2,-4.5) node {$1$};

    % -1 for all three functions
    \draw (-0.7,-0.5) node {$-1$};
    \draw (-0.7,-2.5) node {$-1$};
    \draw (-0.7,-4.5) node {$-1$};

\clip (-4,-6) rectangle (4,4);

\draw[pattern=north west lines, pattern color=blue,draw=white](0,0) arc (0:180:10 and 1);
\draw[draw=blue](0,0) arc (0:180:10 and 1);
\filldraw[white!90!gray] (-3,0.3) circle(5pt);
\draw[color=red] (-3,0.3) node {$+$};


\draw[pattern=north west lines, pattern color=blue,draw=white](0,0) arc (180:0:10 and 1);
\draw[draw=blue](0,0) arc (180:60:10 and 1);
\filldraw[white!90!gray] (3,0.3) circle(5pt);
\draw[color=red] (3,0.3) node {$+$};



\draw[pattern=north west lines, pattern color=blue,draw=white](1,-2) arc (0:180:10 and 1);
\draw[draw=blue](1,-2) arc (0:90:10 and 1);
\draw[pattern=north west lines, pattern color=blue,draw=white](1,-2) arc (180:0:10 and 1);
\draw[draw=blue](1,-2) arc (180:60:10 and 1);
\filldraw[blue] (1,-2) circle(1pt);
\filldraw[white!90!gray] (2.5,-1.8) circle(5pt);
\draw[color=red] (2.5,-1.8) node {$+$};
\filldraw[white!90!gray] (-0.5,-1.8) circle(5pt);
\draw[color=red] (-0.5,-1.8) node {$-$};

\draw[pattern=north west lines, pattern color=blue,draw=white](-1,-4) arc (180:0:10 and 1);
\draw[draw=blue](-1,-4) arc (180:60:10 and 1);
\filldraw[white!90!gray] (2,-3.7) circle(5pt);
\draw[color=red] (2,-3.7) node {$+$};

\draw[pattern=north west lines, pattern color=blue,draw=white](-1,-4) arc (0:180:10 and 1);
\draw[draw=blue](-1,-4) arc (0:180:10 and 1);
\filldraw[white!90!gray] (-3,-3.7) circle(5pt);
\draw[color=red] (-3,-3.7) node {$-$};
\filldraw[blue] (-1,-4) circle(1pt);

  \end{tikzpicture}
\end{center}

У меня получилось 4 промежутка: $(-\infty;-1)$, $(-1;0)$, $(0;1)$, $(1;+\infty)$. С помощью этой картинки очень легко рассуждать. Отдельные выражения ($x^2$, $x-1$, $x+1$) мы перемножаем. Перемножение знаков работает по табличке выше - минус на минус даёт плюс и так далее.

Тогда получается, что на интервале $(-\infty;-1)$ мы умножаем положительный $x^2$ на отрицательный $x-1$ и отрицательный $x+1$, получаем, \textbf{положительный интервал}.

На интервале $(-1;0)$ у нас \textbf{плюс}, \textbf{минус}, \textbf{плюс}, что даёт минус.

На интервале $(0;1)$ у нас опять \textbf{плюс}, \textbf{минус}, \textbf{плюс}, что тоже даёт минус.

Ну и на интервале $(1;+\infty)$ у нас везде \textbf{плюс}, что даёт плюс.

Итого:
\begin{center}
  \begin{tikzpicture}
    % vertical lines
    \draw[->] (0cm,-2) -- (0cm,2cm);
    \draw[dashed] (1,-2) -- (1,1);
    \draw[dashed] (-1,-2) -- (-1,1);

    % axes
    \draw[->] (-5cm,0cm) -- (5cm,0cm) node[right,fill=white] {$x$};

    % points
    \draw (0.2,-0.5) node {$0$};
    \draw (1.2,-0.5) node {$1$};
    \draw (-0.7,-0.5) node {$-1$};

\clip (-4,-2) rectangle (4,4);

\draw[pattern=north west lines, pattern color=blue,draw=white](-1,0) arc (0:180:10 and 1);
\draw[draw=blue](-1,0) arc (0:180:10 and 1);
\filldraw[white!90!gray] (-3,0.3) circle(5pt);
\draw[color=red] (-3,0.3) node {$+$};

\draw[pattern=north west lines, pattern color=blue,draw=white](1,0) arc (180:0:10 and 1);
\draw[draw=blue](1,0) arc (180:60:10 and 1);
\filldraw[white!90!gray] (3,0.3) circle(5pt);
\draw[color=red] (3,0.3) node {$+$};

\draw[pattern=north west lines, pattern color=blue,draw=white](-1,0) arc (180:0:0.5 and 0.4);
\draw[draw=blue] (-1,0) arc (180:0:0.5 and 0.4);
\filldraw[white!90!gray] (-0.5,0.2) circle(5pt);
\draw[color=red] (-0.5,0.2) node {$-$};

\draw[pattern=north west lines, pattern color=blue,draw=white](0,0) arc (180:0:0.5 and 0.4);
\draw[draw=blue] (0,0) arc (180:0:0.5 and 0.4);
\filldraw[white!90!gray] (0.5,0.2) circle(5pt);
\draw[color=red] (0.5,0.2) node {$-$};

\filldraw[blue] (0,0) circle(1.4pt);
\filldraw[blue] (1,0) circle(1pt);
\filldraw[blue] (-1,0) circle(1pt);

  \end{tikzpicture}
\end{center}

Напомню исходное неравенство:
$$f(f(x))\ge{}(f(x))^2$$

Нам нужны положительные промежутки или те точки в которых функция равна нулю. Из картинки видно, что это промежуток $(-\infty;-1]$, промежуток $[1;+\infty)$ и точка ноль.

\section{Ответ}
$$x\in{}(-\infty;-1]\cup{}\{0\}\cup{}[1;+\infty)$$

Важно ${0}$!

\section{Проверка}

\end{document}
